\documentclass[11pt]{article}
\usepackage[a4paper,margin=27mm]{geometry}
\usepackage{amsmath,amssymb,amsthm,hyperref}
\newtheorem{proposition}{Proposition}
\title{The fair reduced signature has exponential Hankel rank}
\author{Algebraic diagnostic; two independent mathematical reviews completed}
\date{30 September 2026}
\begin{document}\maketitle
Fix positive counts $c_1,\ldots,c_n$. Regard the reduced fair signature
as an ordinary noncommutative polynomial by setting
\[
 F(w)=\begin{cases}
 \displaystyle\frac{\prod_{i\in S(w)}c_i}{|w|!},
       &\text{if every letter in $w$ is distinct},\\
 0,&\text{otherwise}.
 \end{cases}
\]
The empty word has coefficient one. Its Hankel matrix is the matrix
$\mathcal H=(F(uv))_{u,v}$ indexed by all finite words.

\begin{proposition}
Over $\mathbb R$, $\operatorname{rank}\mathcal H=2^n$.
\end{proposition}
\begin{proof}
If $u$ repeats a letter, its whole row is zero. Otherwise the row
depends only on its support $S(u)$: the entry is zero if the column
repeats a letter or its support meets $S(u)$, and otherwise depends
only on the union of their supports. There are $2^n$ possible
supports, proving the upper bound.

For each subset $S$ choose one word listing its elements. Restrict
rows to these words and columns to words listing $T^c$, for all
subsets $T$. The entry is nonzero exactly when $S\subseteq T$.
Order both index sets by increasing cardinality, with identical
orders within each cardinality. This matrix is upper triangular,
and its diagonal entries are all
\[
 \frac{\prod_{i=1}^n c_i}{n!}>0.
\]
Thus its rank is $2^n$.
\end{proof}

There is also a direct positive weighted automaton with $2^n$ states.
A state records the set $S$ of letters already read. Reading
$i\notin S$ moves to $S\cup\{i\}$ with weight $c_i/(|S|+1)$;
repeated letters have no transition. Start at the empty set and give
every state terminal weight one. The product of transition weights
is exactly $F(w)$.

\begin{proposition}[A still smaller ordinary completion]
Let $H=\sum_i c_iX_i$ and define the ordinary noncommutative
polynomial $E=\sum_{j=0}^n H^j/j!$, without setting repeated-letter
coefficients to zero. It agrees with all prescribed squarefree fair
coefficients, has nonnegative coefficients, and has Hankel rank $n+1$.
\end{proposition}
\begin{proof}
For a word $u$ of length at most $n$, divide its Hankel row by the
positive product of its letter weights, with repetitions counted.
The resulting row depends only on $|u|$. Longer rows are zero,
so rank is at most $n+1$. Choose rows $X_1^a$ and columns $X_1^b$
for $0\le a,b\le n$. After positive row and column rescaling,
the entries are $1/(a+b)!$ if $a+b\le n$, and zero otherwise.
Reverse the order of the columns. The resulting triangular matrix
has diagonal entries $1/n!$, proving the lower bound.
\end{proof}

\paragraph{Scope.}
These are representations of coefficient data, not constructions
of fair dice. A positive axis word has uncontrolled repeated-letter
coefficients before passage to the reduced algebra; setting these to
zero changes its ordinary Hankel matrix. Consequently no face-count
upper or lower bound follows merely by comparing this rank with the
length of a word. Nor does the result bound the rank of the endpoint
\emph{differential}, whose reduced Lie target has factorial dimension.
The reduced target has an exponential-size prefix--suffix linear
representation, whereas its squarefree constraints admit an ordinary
nonnegative completion of rank only $n+1$. The latter is the truncated
signature of a simultaneous straight path, not of a finite positive
axis word. Thus positivity and the unspecified repeated-letter data
cannot be discarded when seeking a factorial obstruction by
prefix--suffix moment rank.
\end{document}
